% PREPRINT A (user's latest research note): "Residual recovery and unbounded component lifts for Martin norms"
% Notation: [Recovery] and [Check] are companion notes NOT available to us.
% [Check, Theorem 3.1] = "finite local-certificate theorem", stated in Section 1 below.
\documentclass[10pt]{amsart}
\usepackage{amsmath,amssymb,amsthm,mathtools}
\newcommand{\NA}{\operatorname{NA}}
\newcommand{\supp}{\operatorname{supp}}
\newcommand{\sgn}{\operatorname{sign}}
\newcommand{\R}{\mathbb R}
\newcommand{\N}{\mathbb N}
\newcommand{\Cset}{\mathcal C}
\newcommand{\Mset}{\mathcal M}
\title{Residual recovery and unbounded component lifts\\for Mart\'in norms}
\begin{document}
\begin{abstract}
For canonical Mart\'in norms, mate fibers are norm compact at every
normalized functional and depend upper semicontinuously on that
functional. A residual set of non-norm-attaining first rows therefore
admits recovery of every contractive completion into each finite
dimensional Hilbert space. This does not establish density in the
whole operator space. We also recover weighted directions with
unbounded block coefficients. Some resulting non-norm-attaining
operators have no admissible decomposition with bounded first
difference quotients, despite belonging to the closure of
norm-attaining operators. Thus bounded quadratic lifts cannot be
the universal mechanism behind a positive density theorem.
\end{abstract}

\section{Setting and preceding results}
All spaces are real. We use the canonical base on $X=c_0$,
\[
 B_q=B_{c_0}+U(B_H),\qquad
 q^*(a)=\|a\|_1+\|U^*a\|_H,
\]
where $U:H\to c_0$ is compact with dense range. Thus
$\NA(q)=c_{00}$. In the Mart\'in construction \cite{Martin},
$T:\ell_1(\N\times\N)\to X^*$ is injective of norm one,
with range in $Y$ and $Y\cap c_{00}=\{0\}$. Write
\[
 u_{k,m}=\frac{Te_{k,m}}{q^*(Te_{k,m})},\quad
 \Phi_m(k)=2^{-m-k}q^*(Te_{k,m}),\quad
 (R_mx)(k)=m\Phi_m(k)u_{k,m}(x).
\]
Every tail of $(u_{k,m})_k$ is norm dense in $S_{q^*}$.
The space $V_m=(\ell_1,|\cdot|_m)$ has dual norm
$N_m(w)=\|w\|_\infty+\|D_mw\|_2$, where $D_m$ is diagonal
with entries $\Phi_m(k)$. For a nonempty finite block set $I$
or $I=\N$, set
\[
 V=(\bigoplus_{m\in I}V_m)_{\ell_1},\qquad
 Lx=(R_mx)_m,\qquad p(x)=q(x)+\|Lx\|_V.
\]
The bounds $\|R_m\|\leq m2^{-m}$ make $L$ compact and give
$B_{p^*}=B_{q^*}+K$ with $K=L^*(B_{V^*})$ norm compact.
Also $p^{**}(\xi)=q^{**}(\xi)+\|L^{**}\xi\|_V$.

By \cite{Recovery}, $p^{**}$ is strictly convex and every
$f\in S_{p^*}$ has a unique normer $\xi\in S_{p^{**}}$ and
the forced decomposition
\[
 f=a+L^*w=a+\sum_mR_m^*w_m,\qquad
 q^*(a)=1,\quad a(\xi)=q^{**}(\xi),\quad
 w_m=J_{|\cdot|_m}(R_m^{**}\xi).
\]
Here $J_r$ is the norm-one support map for a smooth norm $r$.
For $f\in S_{p^*}$ define
\[
 \Cset(f)=\{g:p^*(f+tg)\leq\sqrt{1+t^2}\ (t\in\R)\},
 \qquad \Mset_f=\bigcup_{n\geq1}n\Cset(f).
\]
The gauge $\gamma_f$ of $\Cset(f)$ satisfies
$\gamma_f(g)^2=\Gamma(f,g):=
\sup_{t\ne0}(p^*(f+tg)^2-1)/t^2$.
The mate inequality is equivalent to contractivity of $(f,g)$;
if $f$ attains its norm, so does $(f,g)$ \cite{KLMW}.

We use the finite local-certificate theorem of \cite{Check}:
if a contractive target has second row
$R_m^*(\omega-dw_m)$, with $\omega$ finitely supported strictly
below the peak of $w_m$ and
$d=\langle D_mw_m,D_m\omega\rangle/\|D_mw_m\|_2$, it is
recoverable whenever its block quadratic coefficient is at most
one. Recovery means membership in
$\overline{\NA((X,p),\ell_2^2)}$.

\section{Global compactness and residual recovery}

Retain the canonical norm $p=q+\|L\cdot\|_V$, with
$q^*(a)=\|a\|_1+\|U^*a\|_H$ and $U,L$ compact.
The graph representation gives
\[
 B_{p^*}=B_{q^*}+K,\qquad K=L^*(B_{V^*}),
 \qquad p^{**}(\xi)=q^{**}(\xi)+\|L^{**}\xi\|_V.
\]
The set $K$ is norm compact. In this section scalar norm
attainment is not required for compactness of the mate fibres.

\begin{theorem}[Global compactness]
\label{thm:gcr-global-compact}
If $f_n\to f$ in $S_{p^*}$ and $g_n\in\Cset(f_n)$, then
$(g_n)$ has a norm-convergent subsequence, and every such limit
belongs to $\Cset(f)$. Consequently every $\Cset(f)$ is norm
compact, and the correspondence $f\mapsto\Cset(f)$ is upper
semicontinuous in norm.
\end{theorem}
\begin{proof}
The vectors $g_n$ are bounded in $\ell_1$. Pass to a weak-star
convergent subsequence $g_n\to g$. The mate inequalities pass
to the limit, so $g\in\Cset(f)$. Choose a normer
$\xi\in S_{p^{**}}$ of $f$ and put $q_0=q^{**}(\xi)>0$.
Then $g(\xi)=0$ and $\|L^{**}\xi\|_V=1-q_0$.

Let $D=\limsup_n\|g_n-g\|_1$, and pass to a subsequence
realizing this limit. Fix $t>0$ and set $s=\sqrt{1+t^2}$.
The dual-ball formula supplies
\[
 f_n+tg_n=a_n+k_n,\qquad q^*(a_n)\le s,
 \qquad k_n\in sK.
\]
After a further subsequence, $k_n\to k$ in norm and
$a_n\to a=f+tg-k$ weak-star. This extraction preserves $D$.
Evaluation at $\xi$ gives
\[
 q^*(a)\ge\frac{a(\xi)}{q_0}
 \ge\frac{1-s(1-q_0)}{q_0}.
\]
Compactness of $U$ gives $U^*a_n\to U^*a$ in norm.
For bounded coordinatewise-convergent sequences in $\ell_1$,
\[
 \|a_n\|_1-\|a_n-a\|_1\longrightarrow\|a\|_1;
\]
this follows by first restricting to finitely many coordinates
and then making the tail of $a$ small. Hence
\[
 tD=\limsup_n\|a_n-a\|_1
 =\limsup_n\big(q^*(a_n)-q^*(a)\big)
 \le\frac{\sqrt{1+t^2}-1}{q_0}.
\]
As $t>0$ was arbitrary, $D=0$. Thus the weak-star subsequence
converges in norm. Compactness of each fibre follows by taking
$f_n=f$. If upper semicontinuity failed, a sequence
$g_n\in\Cset(f_n)$ staying a fixed positive distance from
$\Cset(f)$ would contradict the proved subsequence property.
\end{proof}

\begin{corollary}\label{cor:gcr-meagre-domain}
For every $f\in S_{p^*}$ and every bidual normer $\xi$ of $f$,
$\mathcal M_f$ is a proper, meagre, $K_\sigma$ linear subspace
of $\ker\xi=\{g\in X^*:g(\xi)=0\}$.
\end{corollary}
\begin{proof}
We have $\mathcal M_f=\bigcup_{n\ge1}n\Cset(f)\subseteq\ker\xi$.
Each set in this union is compact and therefore has empty
interior in the infinite-dimensional Banach space $\ker\xi$.
Apply Baire's theorem.
\end{proof}

\begin{corollary}[Lineability without spaceability]
\label{cor:gcr-no-spaceability}
For every $f\in S_{p^*}$, neither $\mathcal M_f$ nor
$\R f+\mathcal M_f$ contains an infinite-dimensional
norm-closed subspace of $X^*$. The gauge $\gamma_f$ makes
$\mathcal M_f$ a Banach space whose inclusion into $X^*$
is compact. At the good scalar norm-attaining supports,
$\mathcal M_f$ is a dense linear subspace of $\ker\xi$
but is not spaceable.
\end{corollary}
\begin{proof}
Both spaces are countable unions of norm-compact sets;
for the second, use the sets $[-n,n]f+n\Cset(f)$.
An infinite-dimensional norm-closed subspace contained in
either would contradict Baire's theorem. Completeness of
the gauge follows from $p^*(g)\le\gamma_f(g)$ and norm
closedness of the gauge balls: an ambient limit of a
gauge-Cauchy sequence is also its gauge limit. Its unit
ball is $\Cset(f)$, so the inclusion is compact.
Density at good scalar supports is the off-peak approximation
result of \cite[Proposition~2.1]{Check}.
\end{proof}

For $d\ge1$, define the joint mate fibre
\[
 \Cset_d(f)=\left\{(g_1,\ldots,g_d):
       f(y)^2+\sum_{j=1}^d g_j(y)^2\le p(y)^2
       \text{ for every }y\in X\right\}.
\]
Thus $\Cset_1(f)=\Cset(f)$. We use any fixed product norm
on $(X^*)^d$ and its associated Hausdorff distance.

\begin{theorem}[Residual uniform recovery]
\label{thm:gcr-residual-recovery}
There is a dense $G_\delta$ set $\Omega\subseteq S_{p^*}$
such that, for every $f\in\Omega$ and every finite $d$,
$\Cset_d$ is Hausdorff continuous at $f$.
In particular, for any sequence
$f_n\in S_{p^*}\cap\NA(p)$ converging to $f$,
\[
 \sup_{G\in\Cset_d(f)}
       \operatorname{dist}(G,\Cset_d(f_n))\longrightarrow0.
\]
Every operator $(f,g_1,\ldots,g_d)$ with
$(g_1,\ldots,g_d)\in\Cset_d(f)$ is consequently a norm
limit of norm-attaining operators into $\ell_2^{d+1}$.
\end{theorem}
\begin{proof}
The set $\Cset_d(f)$ is a closed subset of $\Cset(f)^d$.
Applying Theorem~\ref{thm:gcr-global-compact} to each row shows
that $\Cset_d$ is compact-valued and upper semicontinuous.
Choose a countable dense set $(z_j)$ in $(X^*)^d$.
Each function
\[
 f\longmapsto\operatorname{dist}(z_j,\Cset_d(f))
\]
is finite and lower semicontinuous. By Baire's theorem, all
these functions are continuous on a common dense $G_\delta$
set $\Omega_d$. At a point $f\in\Omega_d$, density of the
$z_j$ gives lower semicontinuity of the fibre: every fixed
$G\in\Cset_d(f)$ is approximated by vectors in nearby fibres.
A finite net of the compact set $\Cset_d(f)$ makes this
approximation uniform over $G$. Together with upper
semicontinuity, this is Hausdorff continuity.

Take $\Omega=\bigcap_{d\ge1}\Omega_d$. Scalar norm-attaining
supports are norm dense by Bishop--Phelps. For $f_n\to f$ as
in the statement and fixed $G\in\Cset_d(f)$, choose
$G_n\in\Cset_d(f_n)$ with $G_n\to G$.
The operator $(f_n,G_n)$ is contractive and has norm one.
At a positive normer of $f_n$, all rows of $G_n$ vanish by
the defining inequality. Hence $(f_n,G_n)$ attains its norm.
\end{proof}

\begin{proposition}[Nontrivial fibres on the residual set]
\label{prop:gcr-nonvacuity}
The set $\Omega$ in Theorem~\ref{thm:gcr-residual-recovery}
contains a dense $G_\delta$ subset $\Omega_0$ consisting of
scalar non-norm-attainers whose mate fibres have
infinite-dimensional linear span.
\end{proposition}
\begin{proof}
For $n\ge1$, put $E_n=\operatorname{span}\{e_1^*,\ldots,e_n^*\}$
and
\[
 F_n=\big[(S_{q^*}\cap E_n)+K\big]\cap S_{p^*}.
\]
Each $F_n$ is norm compact and therefore nowhere dense in
the infinite-dimensional sphere $S_{p^*}$. A support whose
forced base component belongs to $c_{00}$ lies in
$\bigcup_nF_n$. Thus
$\Omega_0=\Omega\setminus\bigcup_nF_n$ is a dense $G_\delta$,
and the forced base component $a$ of every $f\in\Omega_0$
has infinite support. Such $f$ cannot attain its norm,
because scalar norm attainment forces $a\in c_{00}$.

Let $\xi$ be the bidual normer of such an $f$.
The finitely supported vectors $b$ satisfying
\[
 \supp b\subseteq\supp a,\qquad b(\xi)=0
\]
form an infinite-dimensional linear space. For each such $b$,
the function $t\mapsto q^*(a+tb)$ is $C^2$ near zero:
only finitely many nonzero coordinates of $a$ vary, and
$U^*a\ne0$. Its first derivative is
$b(\xi)/q^{**}(\xi)=0$. The local quadratic bound and the
elementary bound away from zero therefore give a finite
$C_b>0$ such that
\[
 q^*(a+tb/C_b)\le\sqrt{1+t^2}\qquad(t\in\R).
\]
Keeping all block supports fixed shows that
$b/C_b\in\Cset(f)$. Hence the linear span of $\Cset(f)$
contains the above infinite-dimensional space.
Scaling any finite collection of independent such vectors
by a sufficiently large common constant gives an independent
tuple in $\Cset_d(f)$.
\end{proof}

For example, if $f\in\Omega_0$, $G\in\Cset_d(f)$ has
independent rows and $0<\rho<1$, then $(f,\rho G)$ is a
non-norm-attaining operator of rank $d+1$ covered by
Theorem~\ref{thm:gcr-residual-recovery}. Its only adjoint
norm-maximizing directions are $\pm e_1$, so norm attainment
would imply that $f$ attains.

Residual recovery does not establish density of all
norm-attaining operators. Upper semicontinuity allows a
fibre at an exceptional support to contain mates that cannot
be approximated from neighbouring fibres. Recovery at those
exceptional supports remains to be proved.

\section{Recovery without bounded component derivatives}

Fix a block $m$, write $R=R_m$, $D=D_m$, and let
$f=a+\sum_lR_l^*w_l\in S_{p^*}$ be the forced decomposition.
Put $w=w_m$, $M=\|w\|_\infty$, $C=\|Dw\|_2$, and
$\lambda_k=m\Phi_m(k)$. Suppose $S\subset\mathbb N$ satisfies
$|w(k)|\le M/2$ for $k\in S$. For a real sequence $\omega$
supported on $S$, assume
\begin{equation}\label{eq:weighted-direction}
 \sum_k\lambda_k\omega(k)^2<\infty.
\end{equation}
The sequence $\omega$ need not be bounded. Nevertheless,
$D\omega\in\ell_2$ and the series
$R^*\omega:=\sum_k m\Phi_m(k)\omega(k)u_{k,m}$ converges
absolutely in $X^*$. Define
\[
 d=\frac{\langle Dw,D\omega\rangle}{C},\qquad
 g=R^*\omega-dR^*w,\qquad
 H=\frac{\|D\omega\|_2^2-d^2}{C}\ge0.
\]

\begin{theorem}[Weighted directions]\label{thm:weighted-recovery}
The direction $g$ has finite mate gauge at $f$, and
\[
 \limsup_{t\to0}\frac{p^*(f+tg)^2-1}{t^2}\le H.
\]
If $(f,g)$ is contractive and $H\le1$, then
$(f,g)\in\overline{\NA((X,p),\ell_2^2)}$.
Consequently some nonzero scalar multiple of every such $g$
has this recovery property.
\end{theorem}
\begin{proof}
Let $\omega_J=P_J\omega$, $d_J=\langle Dw,D\omega_J\rangle/C$,
$g_J=R^*\omega_J-d_JR^*w$, and
$H_J=(\|D\omega_J\|_2^2-d_J^2)/C$. We have $g_J\to g$
in norm and $H_J\to H$. It suffices first to establish the
following estimate, uniformly in $J$, with $J=\infty$ allowed:
\begin{equation}\label{eq:uniform-weighted-expansion}
 p^*(f+tg_J)\le1+\tfrac12H_Jt^2+o(t^2).
\end{equation}

Put $\delta=M/8$ and retain only those coordinates of $\omega_J$
for which $|t\omega(k)|\le\delta$; denote the resulting bounded
sequence by $\omega_{J,t}$. In the dual graph decomposition use
\[
 a+tR^*(\omega_J-\omega_{J,t}),\qquad
 (1-d_Jt)w+t\omega_{J,t}
\]
as the base and $m$th block, keeping all other blocks fixed.
These components sum exactly to $f+tg_J$. The base component
has norm at most
\[
 1+|t|\sum_{|\omega(k)|>\delta/|t|}\lambda_k|\omega(k)|
 \le1+\frac{t^2}{\delta}
       \sum_{|\omega(k)|>\delta/|t|}\lambda_k\omega(k)^2
 =1+o(t^2),
\]
uniformly in $J$.

The numbers $d_J$ are uniformly bounded. For all sufficiently
small $|t|$, the sup norm of the block component is exactly
$(1-d_Jt)M$: the changed coordinates start below $M/2$ and
their perturbations are at most $M/8$. A peak coordinate is
unchanged. The vectors $D\omega_{J,t}$ are uniformly bounded,
and
\[
 \sup_J\|D(\omega_J-\omega_{J,t})\|_2\longrightarrow0.
\]
Moreover, with $d_{J,t}=\langle Dw,D\omega_{J,t}\rangle/C$,
\[
 |t|\,|d_J-d_{J,t}|
 \le\frac{Mt^2}{C\delta}
       \sum_{|\omega(k)|>\delta/|t|}\Phi_m(k)^2\omega(k)^2
 =o(t^2)
\]
uniformly in $J$. Taylor expansion of the Hilbert norm at
$Dw\ne0$ now gives
\[
 N_m((1-d_Jt)w+t\omega_{J,t})
 =1+\tfrac12H_Jt^2+o(t^2)
\]
uniformly in $J$. Since $H_J\ge0$, the dual graph formula
proves \eqref{eq:uniform-weighted-expansion}. The case
$J=\infty$ gives the local assertion. The triangle inequality
away from zero gives a finite full mate gauge.

Suppose next that $(f,g)$ is contractive and $H\le1$, and fix
$0<\rho<1$. The uniform estimate and $H_J\to H$ provide one
neighborhood of zero on which
$p^*(f+\rho t g_J)\le\sqrt{1+t^2}$ for every sufficiently
large $J$. Outside that neighborhood, norm convergence to
$\rho g$ and the strict margin between
$\sqrt{1+t^2}$ and $\sqrt{1+\rho^2t^2}$ give the same
inequality for large $J$. Thus $(f,\rho g_J)$ is contractive.
Its finite off-peak certificate has base curvature zero and
block curvature $\rho^2H_J<1$. The local-certificate recovery
theorem of \cite[Theorem~3.1]{Check} therefore applies. Passing first to $J\to\infty$ and
then to $\rho\uparrow1$ proves the claim. Finally, scaling $g$
by a sufficiently small positive constant makes both its full
mate gauge and its coefficient $H$ at most one.
\end{proof}

\begin{proposition}[Failure of bounded first-order lifting]
\label{prop:no-bounded-lift}
Suppose $a\in c_{00}$ and $\omega$ satisfying
\eqref{eq:weighted-direction} is unbounded and also satisfies
$\sum_k2^{-k}|\omega(k)|<\infty$. Then $g$ has no
representation
\[
 g=b+\sum_lR_l^*v_l,
 \qquad \supp b\subseteq\supp a,\qquad
 (v_l)_l\in V^*.
\]
In particular, no dual graph decomposition of $f+tg$ can have
two-sided norm derivatives at zero as a map into
$X^*\times V^*$, with its component norms bounded by
$1+O(t^2)$. Hence bounded quadratic component lifts do
not exist for every mate.
\end{proposition}
\begin{proof}
The convergent coefficient series defining $g$ belongs to $Y$.
Indeed, its coefficient at $(k,m)$ in the injective map $T$ is
$m2^{-m-k}(\omega(k)-dw(k))$, and these coefficients are
summable by the additional coefficient hypothesis.
Both $g$ and $L^*v$ then belong to $Y$, so the proposed finite
support vector $b$ lies in $Y\cap c_{00}=\{0\}$. Injectivity
of $T$ forces $v_m=\omega-dw$ and $v_l=0$ for $l\ne m$,
contrary to $v\in V^*$.

For the final assertion, differentiate the exact decomposition.
The two-sided base norm bound forces
$\supp b\subseteq\supp a$: the sum of its two directional
derivatives is $2\sum_{k\notin\supp a}|b(k)|$. The preceding
contradiction applies.
\end{proof}

\begin{proposition}[Unavoidable derivative blow-up]
\label{prop:lift-blowup}
Retain the hypotheses of Proposition~\ref{prop:no-bounded-lift},
and suppose that $g$ is a mate and that the normalized base
contact has cube representative $z$ satisfying
$\|z|_{F^c}\|_\infty<1$, where $F=\supp a$.
For every choice of admissible decompositions
\[
 f+tg=A(t)+L^*W(t),\qquad
 q^*(A(t)),\ \|W(t)\|_{V^*}\le\sqrt{1+t^2},
\]
we have $\|W(t)-w\|_{V^*}/|t|\to\infty$ as $t\to0$.
Here $w=(w_l)_l$ is the forced block component of $f$.
\end{proposition}
\begin{proof}
Otherwise choose $t_j\to0$, of one sign, with bounded block
difference quotients. Passing to a weak-star convergent
subsequence gives $(W(t_j)-w)/t_j\to v$ weak-star.
Compactness of $L^*$ gives convergence of their images in norm,
so the base quotients converge in norm to $b=g-L^*v$.
It is enough to treat $t_j>0$; for the other sign replace
$g$ by $-g$.

Evaluation at the bidual normer $\xi$ shows
$b(\xi)\le0$ and $v_l(R_l^{**}\xi)\le0$ for every $l$.
Their absolutely convergent sum is $g(\xi)=0$, so each is
zero. In particular, the upper directional derivative of
$q^*$ at $a$ along $b$ is nonpositive, while $b(\xi)=0$.
Subtracting evaluation at the normalized base contact from
the explicit directional derivative of
$q^*(a)=\|a\|_1+\|U^*a\|_H$ gives
\[
 \sum_{k\notin F}(|b(k)|-z(k)b(k))\le0.
\]
The uniform margin of $z$ forces $b|_{F^c}=0$.
This contradicts Proposition~\ref{prop:no-bounded-lift}.
\end{proof}

\begin{remark}[Non-attaining examples]
The construction applies at non-attaining supports as well.
Fix $a\in c_{00}$, $q^*(a)=1$, set $F=\supp a$, and consider
the open part of the bidual base face
\[
 \left\{z+U\frac{U^*a}{\|U^*a\|_H}:
 z|_F=\sgn a,\quad\|z|_{F^c}\|_\infty<1\right\}.
\]
This is a Baire space in the norm topology of the corresponding
affine $\ell_\infty$ space. For each $N$, the subset where
$|w(k)|<M/2$ for some $k\geq N$ is open and dense.
For density, at a given point $y$ choose $j\notin F$ and
$u\in S_{q^*}$ with $u(y)=0$ and $u(e_j)\ne0$. One may normalize
$e_j^*-y(j)a$. Choose $k_l\to\infty$ with $u_{k_l,m}\to u$ and set
\[
 y_l=y-\frac{u_{k_l,m}(y)}{u_{k_l,m}(e_j)}e_j.
\]
Then $y_l\to y$ within the open face, and $u_{k_l,m}(y_l)=0$,
so the selected block support coordinate is zero.
Openness follows from continuity of the finite weighted support
data. Intersecting these open dense sets gives a residual subset
with infinitely many half-peak indices. This subset contains points
outside $c_0$, since the intersection of $c_0$ with the affine space
is closed and has empty interior. Choose such a point and
normalize it in $p^{**}$. Its support $f$ is not norm-attaining,
by strict convexity of $p^{**}$, and it has infinitely many
indices $k_j$ with $|w(k_j)|<M/2$.

Set $\omega(k_j)=j$ and zero elsewhere. Since $k_j\ge j$,
both \eqref{eq:weighted-direction} and summability of the
$T$-coefficients hold. A sufficiently small nonzero multiple
$g_0$ of the resulting $g$ is a mate with coefficient at most
one. Theorem~\ref{thm:weighted-recovery} recovers $(f,g_0)$,
whereas Proposition~\ref{prop:no-bounded-lift} excludes every
bounded first-order component lift. For $0<\rho<1$,
$(f,\rho g_0)$ itself is non-norm-attaining: the strict damping
forces equality in the original contraction only when the
second coordinate is zero and $f$ attains its norm. Thus there
are non-norm-attaining operators in the closure of norm-attainers
which are inaccessible to any bounded quadratic-lift argument.
Proposition~\ref{prop:lift-blowup} also shows that every admissible
decomposition of these damped targets has unbounded first
difference quotients.
\end{remark}

\section{Consequences for the remaining density problem}
The residual recovery theorem concerns the first-row sphere.
It is not a statement that recoverable operators form a residual
subset of the whole operator space. In particular, a compact
upper semicontinuous correspondence may have additional fiber
directions at discontinuity points; density of continuity points
does not imply density of their graph in the full graph.

The weighted construction supplies a separate improvement.
Bounded component derivatives are not necessary even for
recoverable non-norm-attaining targets. The relevant replacement
is a truncation whose discarded pullback costs $o(t^2)$ uniformly
over the finite approximants. Proving such a representation for
arbitrary mates, or recovering the additional fibers at all
discontinuity supports by another method, remains open here.
Neither compactness nor the derivative blow-up proved above is
a counterexample to operator density.

\begin{thebibliography}{9}
\bibitem{KLMW}
V.~Kadets, G.~L\'opez, M.~Mart\'in and D.~Werner,
\emph{Norm attaining operators of finite rank}, arXiv:1905.08272.
\bibitem{Martin}
M.~Mart\'in,
\emph{A Banach space whose set of norm-attaining functionals is
algebraically trivial}, J. Funct. Anal. \textbf{288} (2025), 110815.
\bibitem{Recovery}
\emph{Recovery of mates on Mart\'in norms and generic second-order
curvature}, accompanying research note, October 8, 2026. [NOT AVAILABLE]
\bibitem{Check}
\emph{Checking the positive density route for Mart\'in norms},
accompanying research note, October 8, 2026. [NOT AVAILABLE]
\end{thebibliography}
\end{document}
